\section{Direct contraction by conditioning}\label{sec:conditioning}
Let $G=\bm y-\bm u$, $\bm u=K^{-1}\bm1$, $\alpha=e^{2\pi\mathrm i/K}$ and $\rho_j=1-p+p\alpha^j$. Fix $s_0$ and the query. For $S\subseteq[n]$, let $\mathcal F_S=\sigma(w_t:t\in S)$.

\begin{lemma}[Conditional target contraction]\label{lem:conditional-target}
For every $S\subseteq[n]$,
\[
 \|\mathbb E[G\mid\mathcal F_S]\|^2
 =\frac1K\sum_{j=1}^{K-1}|\rho_j|^{2(n-|S|)}
 \le (1-K^{-1})\lambda^{2(n-|S|)}.
\]
For every square-integrable $\mathcal F_S$-measurable vector $H$, and every matrix-valued such $H$ with compatible dimensions,
\[
 |\mathbb E\langle H,G\rangle|\le\lambda^{n-|S|}\|H\|_{L^2},\qquad
 \|\mathbb E[H^\top G]\|\le\lambda^{n-|S|}\|H\|_{L^2(F)}.
\]
\end{lemma}
\begin{proof}
The nonconstant Fourier coefficient of the count distribution at frequency $j$ is the conditional expectation of $\alpha^{js_n}$. Independence gives
\[
 \mathbb E[\alpha^{js_n}\mid\mathcal F_S]
 =\alpha^{js_0}\rho_j^{n-|S|}\prod_{t\in S}\alpha^{jb_t}.
\]
Every factor in the product has modulus one. Parseval for the $K$-point discrete Fourier transform, with its $1/K$ normalization on probability vectors, gives the asserted identity and bound. Since $H$ is measurable on $\mathcal F_S$, the tower property gives $\mathbb E\langle H,G\rangle=\mathbb E\langle H,\mathbb E[G\mid\mathcal F_S]\rangle$. Cauchy--Schwarz proves the scalar bound. Testing the matrix expression against an arbitrary unit vector proves the vector bound. The factor $\sqrt{1-1/K}\le1$ was dropped.
\end{proof}

\paragraph{Which objects require conditioning.} The argument applies to a network term supported on $S$, even when that term has nonzero conditional means. It therefore permits summing the ordinary attention-expansion terms directly. No decomposition by orthogonal interaction degree is used in the outcome proof below.

\section{Direct local credit and the outcome plateau}\label{sec:direct-outcome}
For fixed parameters, introduce an auxiliary variable $u$: multiply every content attention score, in both blocks, by $u$, leaving any positional relative bias unchanged. The second-block residuals and hence its content scores depend on $u$. This is the composed network, rather than an expansion with frozen second-block scores. Define
\begin{equation}
 z_\theta=\sum_{k=0}^N z^{[k]}+R_N,\qquad
 z^{[k]}=\frac1{k!}\frac{d^kz(u)}{du^k}\bigg|_{u=0},\qquad
 R_N=\int_0^1\frac{(1-u)^N}{N!}\frac{d^{N+1}z(u)}{du^{N+1}}\,du.
 \label{eq:score-split}
\end{equation}
For credit Jacobians let $J^{[k]}=\partial_\theta z^{[k]}$ and $J_R=\partial_\theta R_N$. The auxiliary Taylor order is not the parameter-polynomial degree.

\begin{lemma}[Support and absolute coefficient bounds]\label{lem:direct-support}
On the doubled initialization ball, with relative bias fixed at zero, write every positive-order coefficient as a finite sum of terms $z^{[k]}=\sum_\nu H_{k,\nu}$, each depending on at most $2k+1$ letters. There are constants $C,B$ independent of $n$, widths and $k$ such that for $k\ge1$,
\[
 \sum_\nu\sup\|H_{k,\nu}\|\le C\varepsilon^3(B\varepsilon^4)^k,\qquad
 \sum_\nu\sup\|\partial_\theta H_{k,\nu}\|_F\le C\varepsilon^2(B\varepsilon^4)^k.
\]
The remainder satisfies $\sup\|R_N\|\le C\varepsilon^3(B\varepsilon^4)^{N+1}$ and $\sup\|J_R\|_F\le C\varepsilon^2(B\varepsilon^4)^{N+1}$. At the answer position the zero-order coefficient is a fixed-input term of size $C\varepsilon^2$ and tangent size $C\varepsilon$, plus a sum of one-letter terms with absolute sizes at most $C\varepsilon^3$ and tangent sizes at most $C\varepsilon^2$.

The logit assertions hold uniformly for arbitrary relative bias $\beta$, provided the other parameters stay in the doubled ball.
\end{lemma}
\begin{proof}
The complete product, softmax, two-block and marked-site calculations are given in \cref{prop:network}. Its absolute coefficient bounds give the first assertion before any cancellation; its input-count calculation gives support at most $2k+1$. For a map into $\mathbb R^K$, its parameter Jacobian has Frobenius norm at most $\sqrt K$ times its operator norm. Thus the directional tangent bounds there give the stated Frobenius bounds, with a fixed $\sqrt K$ absorbed in $C$. The integral in \eqref{eq:score-split} has normalized weight exactly matching the next normalized derivative, proving the remainder estimate.

At zero scores the attention weights are input independent. The direct answer-position embedding and every MLP offset depend on no random letters (conditional on the start and query). All transported letter terms contain the readout, a value matrix and a letter embedding. Their sizes are $O(\varepsilon^3)$, their tangents $O(\varepsilon^2)$, and their coefficients sum absolutely to a constant because all attention rows are probability vectors. This proves the zero-order assertion.

For arbitrary $\beta$, replace the zero-score row $m_i^{-1}\bm1$ by $\operatorname{softmax}(\beta_{i-\cdot})$. Every such row has nonnegative entries summing to one. The derivative identity and normalized softmax derivative estimate in \eqref{eq:softmax-avg} hold at every real score vector. At $u=0$ the coefficients are weighted averages with those fixed probabilities; the absolute sum estimate is unchanged. Biases add no input sites and do not differentiate with respect to $u$, so the support induction is unchanged as well. The residual bounds imply $|(W_{\rm in}h)_m|\le C\varepsilon^2<|b_m|$, independently of $\beta$, keeping the activation pattern fixed. Only the marked-site $1/m_i$ estimate in \cref{prop:network}(iv) requires $\beta=0$; it is not used for the logit assertions here.
\end{proof}

\begin{lemma}[Direct credit summation]\label{lem:score-sum}
Let $0<\lambda<1$, $N=\lceil(n-1)/2\rceil$ and $B\varepsilon^4\le\lambda^2/2$. Under the local assumptions,
\begin{align*}
 |\outcredit(\theta)|&\le C(\varepsilon^2\lambda^n+\varepsilon^3\lambda^{n-1}+\varepsilon^7\lambda^{n-3}),\\
 \|\nabla_\theta\outcredit(\theta)\|&\le C(\varepsilon\lambda^n+\varepsilon^2\lambda^{n-1}+\varepsilon^6\lambda^{n-3}).
\end{align*}
The first bound is uniform in arbitrary relative bias. In the second bound the bias is fixed at zero and excluded from the parameter Jacobian. If also $\varepsilon\le\lambda$, these bounds reduce to
\begin{equation}
 |\outcredit(\theta)|\le C\varepsilon^2\lambda^n,\qquad
 \|\nabla_\theta\outcredit(\theta)\|\le C\varepsilon\lambda^n.
 \label{eq:thm1-full}
\end{equation}
They remain valid after averaging over the start count.
\end{lemma}
\begin{proof}
Apply \cref{lem:conditional-target} directly to each coefficient term in \cref{lem:direct-support}. The zero-order bounds give $C(\varepsilon^2\lambda^n+\varepsilon^3\lambda^{n-1})$ for the logit credit and $C(\varepsilon\lambda^n+\varepsilon^2\lambda^{n-1})$ for its gradient. For $k\ge1$, every support size is at most $2k+1$, so
\[
 |\mathbb E\langle z^{[k]},G\rangle|
 \le C\varepsilon^3(B\varepsilon^4)^k\lambda^{n-2k-1}
 = CB\varepsilon^7\lambda^{n-3}
 \left(\frac{B\varepsilon^4}{\lambda^2}\right)^{k-1}.
\]
If $2k+1>n$, the displayed power of $\lambda$ is larger than one and still gives a valid upper bound. For the remainder use $\|G\|\le1$, hence
\[
 |\mathbb E\langle R_N,G\rangle|
 \le C\varepsilon^3(B\varepsilon^4)^{N+1}
 =CB\varepsilon^7\lambda^{2N}
 \left(\frac{B\varepsilon^4}{\lambda^2}\right)^N
 \le CB\varepsilon^7\lambda^{n-3}\,2^{-N},
\]
since $2N\ge n-1$. Summing a geometric series proves the first bound. Apply the same conditional estimate to the Jacobian terms, with one fewer power of $\varepsilon$, for the second bound. If $\varepsilon\le\lambda$ and $\varepsilon\le1$, then $\varepsilon/\lambda\le1$ and $\varepsilon^5/\lambda^3\le\varepsilon^2\le1$, which proves \eqref{eq:thm1-full}. All bounds were uniform in the fixed start, so averaging preserves them.
\end{proof}

\begin{lemma}[Uniform loss floor and smoothness]\label{lem:bias-uniform}
On the set $\mathcal B^{\rm full}_\varepsilon=\{(\vartheta,\beta):\|\vartheta-\vartheta_0\|\le2r_0\varepsilon\}$, with $\beta$ unrestricted, if $\varepsilon\le c\lambda$ then
\[
 |\outcredit|\le C_2\varepsilon^2\lambda^n.
\]
The loss has gradient norm at most $C_1\varepsilon$ and directional second derivatives at most $C_1$, for the full parameter norm including $\beta$.
\end{lemma}
\begin{proof}
Choose $c$ small enough that $\varepsilon\le\lambda$ and $B\varepsilon^4\le\lambda^2/2$. The credit estimate follows from \cref{lem:score-sum}. For smoothness, a unit perturbation of the relative bias changes a row of scores by a vector of supremum norm at most one. By \eqref{eq:softmax-avg}, its attention row has bounded first and second derivatives in $\ell_1$, uniformly in its number of visible sites and in $\beta$. All parameter matrices and input embeddings have directional profiles $(C\varepsilon,C,0)$. A bias vector has norm at most $C$ and profile $(C,1,0)$; it enters through $o=W_{\rm out}D b$, whose profile is $(C\varepsilon,C,C)$. Products have profiles
\[
 (A,A',A'')(B,B',B'')=(AB,A'B+AB',A''B+2A'B'+AB'').
\]
Induct over the two blocks. Residuals have profile $(C\varepsilon,C,C)$; projected queries, keys and values have profile $(C\varepsilon^2,C\varepsilon,C)$; content scores have profile $(C\varepsilon^4,C\varepsilon^3,C\varepsilon^2)$. Including the relative bias leaves attention rows with profile $(1,C,C)$ in $\ell_1$. A head output has profile $(C\varepsilon^2,C\varepsilon,C)$ and the final logits have profile $(C\varepsilon^2,C\varepsilon,C)$. The fixed activation pattern makes every intervening MLP affine. For a single-example cross-entropy, $p=\operatorname{softmax}(z)$,
\[
 \ell'=\langle p-y,z'\rangle,\qquad
 \ell''=z'^\top(\operatorname{diag}p-pp^\top)z'+\langle p-y,z''\rangle.
\]
Since $\|p-y\|\le\sqrt2$ and $\|\operatorname{diag}p-pp^\top\|\le1/2$, these are bounded by $C\varepsilon$ and $C$. Taking expectations and testing arbitrary unit directions proves the assertions. Every estimate used operator norms or probability-row $\ell_1$ norms, so the constants have the claimed independence.
\end{proof}
\label{lem:end-smooth}

\paragraph{Loss identity.}\label{app:end}
For $p_\theta=\operatorname{softmax}(z_\theta)$, direct algebra gives
\begin{equation}
 L_{\rm out}(\theta)=\log K+\mathbb E D_{\rm KL}(u\|p_\theta)-\outcredit(\theta).
 \label{eq:lo-identity}
\end{equation}
Indeed $D_{\rm KL}(u\|p)=\log\sum_a e^{z_a}-\log K-\langle u,z\rangle$ and the cross-entropy is $\log\sum_a e^{z_a}-\langle y,z\rangle$.

\begin{proof}[Complete proof of \cref{thm:outcome-end}]
Set $B_n=C_2\varepsilon^2\lambda^n$. On the doubled local set, \eqref{eq:lo-identity} and nonnegativity of KL give
\begin{equation}L_{\rm out}(\theta)\ge\log K-B_n.\label{eq:lo-floor}\end{equation}
At zero initial readout, $L_{\rm out}(\theta_0)=\log K$. Let $\tau$ be the first exit time from the closed full parameter ball of radius $r_0\varepsilon$, including the relative bias. For piecewise constant block projections (with finitely many switches on each finite time interval), local smoothness gives a unique trajectory between switches and a continuous trajectory across them. Almost everywhere before exit,
\[
 \frac{dL_{\rm out}}{dt}=-\|P(t)\nabla L_{\rm out}\|^2=-\|\dot\theta\|^2.
\]
Consequently, for $t\le\tau$,
\[
 \int_0^t\|\dot\theta(s)\|^2ds\le B_n,\qquad
 \|\theta(t)-\theta_0\|^2\le t\int_0^t\|\dot\theta(s)\|^2ds\le tB_n.
\]
If exit occurs, continuity gives $r_0^2\varepsilon^2\le\tau B_n$, hence $\tau\ge r_0^2 C_2^{-1}\lambda^{-n}$. If it never occurs, this lower bound is automatic. Taking $c\le r_0^2/(2C_2)$ proves locality on the claimed interval and proves the loss floor there. The argument also holds for measurable block projections using an absolutely continuous flow, since the same energy identity holds almost everywhere; thus permitting accumulation of switching times does not alter the bound.

Finally descent and the identity give $\mathbb E D_{\rm KL}(u\|p_t)\le\outcredit(\theta(t))\le B_n$. For every example, $p_{t,y}-1/K\le\operatorname{TV}(p_t,u)$. Pinsker's inequality and Jensen therefore give
\[
 \mathbb E[p_{t,y}]\le K^{-1}+\sqrt{B_n/2}
 \le K^{-1}+C\varepsilon\lambda^{n/2}.
\]
For completeness, Pinsker's inequality here follows by grouping the coordinates where $u_a\ge p_a$, applying the log-sum inequality to the two groups, and using $D_{\rm KL}((x,1-x)\|(q,1-q))\ge2(x-q)^2$. The last inequality follows by differentiating twice in $x$: its second derivative is $1/[x(1-x)]\ge4$ and it and its first derivative vanish at $x=q$. All time constants are independent of the block schedule.
\end{proof}
